\chapter{Troubleshooting}\label{ch:troubleshooting}

This chapter collects problems that users run into most often, explains their causes and shows
how to fix them. It is organised by symptom.

\section{Reading error messages}\label{sec:ts-errors}

Most failures surface as \code{nf.GeometryError}. The message starts with the name of the
operation and states the violated condition, often with the offending values. The following
script provokes a few typical mistakes:

\example{18_common_errors.py}

\begin{lstlisting}[style=text]
knot vector: GeometryError: need at least degree+1 control points (got 2)
zero axis: GeometryError: solid primitive: axis must not be zero
offset > radius: GeometryError: offset_curve: |distance| = 1.500000 reaches the radius of curvature 1.000000 at u = 0.000000 on the side of the centre of curvature (the offset would have a cusp)
2D point: TypeError: make_line(): incompatible function arguments. The following argument types are supported:
angle in degrees: GeometryError: revolution angle must be in (0, 2*pi]
angle in radians: ok
open shell closed: False  free edges: 4
boolean with open shell: GeometryError: boolean: both shells must be closed
\end{lstlisting}

\begin{description}[style=nextline,leftmargin=1.5em]
  \item[\code{TypeError: ... incompatible function arguments}] An argument has the wrong type
    or shape. The full message lists every accepted signature (overloads). Common causes:
    points with two coordinates (give $z=0$), a single number where a sequence is expected,
    a \code{NurbsSurface} where a \code{TrimmedSurface} or \code{Shell} is needed, or an enum
    value given as a string (\code{nf.BooleanOp.Union}, not \code{"union"}).
  \item[Errors about knot vectors] The knot vector must have \code{n_ctrl + degree + 1}
    entries, be non-decreasing, and have interior multiplicities of at most \code{degree}.
    Use \apiref{clamped_uniform_knots} to create a valid one and \apiref{validate_knots} to
    check your own.
  \item[\code{... must be in (0, 2*pi]}] Angles are radians everywhere in the Python and C++ API.
    Only the LLM tools take degrees (\cref{ch:llm}).
\end{description}

\section{Results that look wrong}\label{sec:ts-wrong}

\paragraph{A method call ``does nothing''.}
Geometry objects are immutable (\cref{sec:concepts-immutability}). \code{c.insert_knot(0.5)},
\code{shell.oriented()} and \code{s.transformed(T)} return new objects. Assign the result:
\code{c = c.insert_knot(0.5)}.

\paragraph{A circle starts at an unexpected point.}
\code{make_circle} chooses the start direction perpendicular to the normal itself (for the
normal $+z$ the $+y$ axis). If the start point matters, for example for a seam position or for
lofting circles, use \apiref{make_arc} with explicit axes or \apiref{make_circle_3pt}.

\paragraph{Parameters are not proportional to length.}
$u=0.5$ is not the midpoint by length in general, not even for circles made of several rational
arcs. Use \apiref{parameter_at_length} to find a point at a given arc length.

\paragraph{Offsets or fillets go to the wrong side.}
Curve offsets go to the \emph{left} of the curve direction seen from \code{plane_normal} for a
positive distance. Surface offsets and surface fillet sides refer to the normal
$\mathbf S_u\times\mathbf S_v$. Check the direction with \code{surface.normal(u, v)}, reverse the
sign of the distance, or reverse or transpose the input.

\paragraph{Normals of a shell point inwards; the volume is negative.}
Shells sewn from faces are not automatically oriented. Call \apiref{Shell.oriented} (it orients
every connected component outwards) and check \code{shell.is_consistently_oriented}. Solid
primitives and Boolean results are already oriented.

\paragraph{Printed values are zero.}
The text representation of \code{Vec3} and \code{Tolerance} shows six decimal places, so
$10^{-10}$ prints as \code{0.000000}. Print the attributes with a format string.

\section{Shells and Boolean operations}\label{sec:ts-booleans}

\paragraph{\code{boolean: both shells must be closed}.}
Booleans, volumes and point classification need closed shells. Check
\code{shell.is_closed} and the edges listed by \apiref{Shell.free_edges}. Common causes are faces
that do not meet within the sewing tolerance (raise \code{Shell.sew(faces, tolerance=...)} for
models in large units, or fix the geometry), a missing face, or trim loops that do not match the
neighbouring face exactly.

\paragraph{Degenerate configurations.}
Boolean operations are most robust when the operands intersect \emph{transversally}. The
algorithm detects coincident faces and cuts them along each other's boundaries. It treats
tangential contacts and overlaps thinner than the tolerance as contacts, so two boxes sharing a
face, or a tool flush with a face, are handled. Problems can still arise when
\begin{itemize}
  \item coincident faces have a common part that crosses the seam of a closed surface (a known
    limitation),
  \item an intersection curve passes exactly through the apex of a cone,
  \item faces are tangent along a curve that is not a face boundary.
\end{itemize}
Curves through the poles of a sphere and along the seam of a closed surface are handled (a
sphere minus a box with its corner at the centre is closed on all platforms; up to version
0.9.0 it was open with GCC and MSVC, issue \#63). If a result is open nevertheless, avoid the
exact degeneracy: extend tools slightly beyond the faces they cut (make a drill longer than
the plate is thick) or move a cut a small distance.

\begin{tip}[Checking Boolean results]
After a Boolean, check \code{result.is_closed}, compare \code{result.volume()} with a rough
expectation (for instance the operands' volumes), and look at \code{len(result.faces)}. The LLM
interface does exactly this in its verification rounds (\cref{ch:llm}).
\end{tip}

\paragraph{Booleans are slow.}
Run time grows with the number of face pairs and the complexity of the intersection curves.
Rational faces of high degree from sweeps or fits cost more than primitives. Build solids
from primitives where possible, and combine many tools into one before subtracting them from a
part.

\section{Tolerances}\label{sec:ts-tolerances}

\paragraph{The requested tolerance cannot be met.}
For example \code{sag_surface: tolerance not reached with 256 spans}. Approximating
operations refine until the tolerance is met, but they stop at a limit
(\code{max_sections=1025} for sweeps and surface fillets, \code{max_spans=256} for
\code{sag_surface}). Loosen the tolerance, or simplify the input. Very tight tolerances
($<10^{-9}$ relative to the model size) approach floating-point resolution and may never be
met.

\paragraph{Models in very large or very small units.}
The default sewing and Boolean tolerances ($10^{-6}$) and most fit tolerances ($10^{-7}$) are
absolute, in model units. They suit models of size 1 to 1000. For a part in metres of size 0.01,
or in micrometres of size $10^5$, scale the tolerances, or change the unit of the model
(\cref{sec:concepts-units}).

\paragraph{Results are bigger than necessary.}
Approximations get more knots the tighter the tolerance. A surface fitted to $10^{-9}$ can have
many times the control points of one fitted to $10^{-6}$. This slows down intersections,
Booleans and meshing. Request only the accuracy you need.

\section{Python and NumPy}\label{sec:ts-python}

\paragraph{Point lists and grids.} Functions for many points take NumPy arrays: point lists of
shape $(n,3)$ and grids of shape $(n_u,n_v,3)$. An array of the wrong shape, e.g.\ $(3,n)$,
raises an error. Transpose it with \code{pts.T}. Single points accept any sequence of three numbers,
NumPy rows included.

\paragraph{\code{ImportError: ... _core}.} The compiled extension does not match the interpreter,
or the source tree shadows the installed package because Python runs from the repository root
without an editable installation. Reinstall with the interpreter you use (\code{python -m pip
install ...}), or run from another directory.

\paragraph{Objects cannot be pickled.} Up to version 0.9.0 the geometry objects could not be
pickled; pass geometry between processes as STEP text (\apiref{to_step}, \apiref{from_step})
there. Newer versions pickle all geometry types (\cref{sec:concepts-threads}).

\section{gmsh and meshing}\label{sec:ts-gmsh}

\paragraph{\code{OSError} on \code{import gmsh} (Linux).} The gmsh wheel needs X11 and OpenGL
runtime libraries even when used headless (\cref{sec:install-gmsh-linux}). Install them, or
catch \code{(ImportError, OSError)} if meshing is optional.

\paragraph{\code{MeshTopologyError: surface mesh is not watertight}.} The surface mesh has open
or non-manifold edges. Mesh a \emph{closed} shell added with \code{add_shell}. Surfaces added one
by one with \code{add_surface} do not share their edges. Check the shell with
\code{shell.is_closed} first.

\paragraph{Faces missing in gmsh, wrong volumes.} Shells are transferred to gmsh through the
native STEP writer and OpenCASCADE's STEP reader, so that OCC builds the seam edges of trimmed
closed faces itself. Versions before 0.9.0 transferred faces one by one, which lost such faces
after Booleans. Update if you see this with an older version.

\paragraph{The mesh is too coarse or too fine near small features.} Use
\apiref{mesh.GmshSession.refine} on the face labels of the features, with a \code{distance} for a
smooth transition, and \code{curvature} in \code{mesh()} for curved faces (\cref{ch:meshing}).

\section{Platform-specific issues}\label{sec:ts-platform}

\paragraph{macOS: \code{Nurbforge needs macOS >= 13.3}.} CMake stops if the deployment target is
below 13.3, because the STEP writer needs \code{std::to_chars} for floating-point numbers, which
Apple's C++ library provides only from that version. Set
\code{MACOSX_DEPLOYMENT_TARGET=14.0} (as the release wheels do) or leave it unset to build for the
running system.

\paragraph{Locale and decimal separators.} STEP files always use a decimal point. The writer
and reader do not depend on the C locale. Reading works even in a process whose
\code{LC_NUMERIC} uses a decimal comma, as with German or French locale settings. On macOS the
reader parses numbers with a locale-independent \code{strtod_l}, and the test suite checks this
with a \code{de_DE} locale.

\paragraph{Windows: \code{Python was not found}, \code{NMake Makefiles}.} See
\cref{sec:install-windows}.

\paragraph{The C++ package is not found: \code{Threads::Threads ... target was not found}.}
This happens with packages installed from version 0.9.0 or older: call
\code{find_package(Threads REQUIRED)} before \code{find_package(Nurbforge)}, or install a newer
version, whose package configuration finds the dependency itself (\cref{sec:cpp-cmake}).

\section{The LLM interface}\label{sec:ts-llm}

\paragraph{\code{OllamaError: cannot reach Ollama at http://localhost:11434}.} The Ollama server
is not running or listens elsewhere. Start it (\code{ollama serve}, or the desktop app) and check
\code{OLLAMA_HOST}. A host without scheme, such as \code{127.0.0.1:11434}, is accepted;
\code{http://} is added.

\paragraph{The model does not call tools, or answers in prose only.} The model must support tool
calling (\code{qwen2.5}, \code{llama3.1}, \code{mistral-nemo} and others do). Small models
handle multi-step requests poorly. Use a larger model, split the request, or state dimensions
explicitly.

\paragraph{The result has wrong dimensions.} Keep the verification rounds enabled
(\code{verify_rounds=1}, the default). For dimensions that are only fixed indirectly (``a
cylinder with lateral area 50''), ask the model to use the \toolref{solve} tool
(\cref{ch:llm}).

\section{Performance tips}\label{sec:ts-performance}

\begin{itemize}
  \item Evaluate in bulk: \code{curve.sample(n)} and \code{surface.sample(nu, nv)} return NumPy
    arrays and avoid one Python call per point.
  \item Prefer exact constructions (primitives, revolutions, extrusions) to fitted ones. They are
    smaller and faster in all later operations.
  \item Use the loosest tolerance that meets your needs; size and run time of approximations grow
    with the required accuracy.
  \item Reuse objects: they are immutable, so a fitted surface can be shared by many shells and
    transforms without copies.
  \item Parallelism: the surface/surface intersection is multi-threaded internally.
    The expensive operations release Python's interpreter lock, so variants of a parameter
    study can run in threads; processes work as well (\cref{sec:concepts-threads}).
  \item Mass properties and areas take \code{rel_tol}; the default $10^{-8}$ is more than enough
    for most engineering purposes, and $10^{-6}$ is noticeably faster on complex shells.
\end{itemize}
